\documentclass[12pt]{article}
\usepackage[utf8]{inputenc}
\usepackage{geometry}
\usepackage{hyperref}
\usepackage{enumitem}
\usepackage{tabularx}
\usepackage{colortbl}
\usepackage{xcolor}
\usepackage{booktabs}
\usepackage{titlesec}

% Page setup
\geometry{a4paper, margin=1in}
\hypersetup{
	colorlinks=true,
	linkcolor=blue,
	urlcolor=blue,
	citecolor=blue,
}

% Section formatting
\titleformat{\section}{\large\bfseries\color{darkblue}}{}{0em}{}
\titleformat{\subsection}{\normalsize\bfseries\color{blue}}{}{0em}{}

% Custom colors
\definecolor{darkblue}{rgb}{0.1,0.2,0.5}
\definecolor{lightgray}{gray}{0.95}

\begin{document}
	
	% Title
	\begin{center}
		{\Huge\bfseries Training Program for Big Data \& Distributed Systems} \\[0.5cm]
		{\large Bridge Role — Supporting Until Permanent Position} \\[0.3cm]
		{\textit{Structured, Intensive, Career-Level Training}} \\[1cm]
	\end{center}
	
	\begin{quote}
		\emph{This is a bridge role intended to support you until you secure a permanent position. The training path is structured and intensive. Realistically, it takes 6 months to 2 years to gain full readiness for the role, depending on consistency and depth of practice.}
		
		The courses must be taken in the order listed below, as each builds on the previous one. Training will focus heavily on \textbf{Python}, \textbf{distributed systems}, \textbf{testing}, and \textbf{big-data computation}.
		
		Let me know once you start, and we will review progress weekly.
		
		Best regards, \\
		Satyadhar
	\end{quote}
	
	
	% ====================================================
	% COURSE LINKS (ADDED — NO REMOVALS)
	% ====================================================
	\section*{Complete Course Links (Mandatory Sequence)}
	
	\begin{enumerate}[label=\textbf{Course \arabic*:}]
		\item Remote \& Independent Working in Managerless Environments  
		\href{https://www.udemy.com/course/remote-independent-working-in-a-managerless-environment-101/}
		{https://www.udemy.com/course/remote-independent-working-in-a-managerless-environment-101/}
		
		\item Big Data Infrastructure Administration — Ticket-Based Learning  
		\href{https://www.udemy.com/course/big-data-infra-admin-101-ticket-based-learning/}
		{https://www.udemy.com/course/big-data-infra-admin-101-ticket-based-learning/}
		
		\item Python Full-Stack Computational Framework (Beginner)  
		\href{https://www.udemy.com/course/python-full-stack-computational-framework-for-begineers-101/}
		{https://www.udemy.com/course/python-full-stack-computational-framework-for-begineers-101/}
		
		\item Full-Spectrum Comprehensive Testing \& Error Reading  
		\href{https://www.udemy.com/course/full-spectrum-comprehensive-testing-error-reading-101/}
		{https://www.udemy.com/course/full-spectrum-comprehensive-testing-error-reading-101/}
		
		\item Python Big Data Refresher — Intermediate (Quiz Based)  
		\href{https://www.udemy.com/course/refresher-python-big-data-intermediate-quiz-based-course-102/}
		{https://www.udemy.com/course/refresher-python-big-data-intermediate-quiz-based-course-102/}
		
		\item Python Full-Stack Backend \& ML Engines  
		\href{https://www.udemy.com/course/python-full-stack-and-backend-engines-for-mc-ml-engines-102/}
		{https://www.udemy.com/course/python-full-stack-and-backend-engines-for-mc-ml-engines-102/}
		
		\item Writing Tests for Simulation Engineering \& Code Conversion  
		\href{https://www.udemy.com/course/writing-tests-for-simeng-python-code-conversion-concepts-101/}
		{https://www.udemy.com/course/writing-tests-for-simeng-python-code-conversion-concepts-101/}
		
		\item Full-Stack Computational Framework for Big Data Simulations  
		\href{https://www.udemy.com/course/full-stack-computational-framework-for-big-data-simulations/}
		{https://www.udemy.com/course/full-stack-computational-framework-for-big-data-simulations/}
		
		\item Converting Standalone Python Code to Distributed Analytics Code  
		\href{https://www.udemy.com/course/converting-standalone-code-to-distributed-code-for-analytics/}
		{https://www.udemy.com/course/converting-standalone-code-to-distributed-code-for-analytics/}
	\end{enumerate}
	
	\newpage
	
	% ====================================================
	% ASSUMPTIONS (ORIGINAL PRESERVED)
	% ====================================================
	\section*{Training Assumptions}
	
	\begin{itemize}
		\item Each course requires approximately 2--4 days
		\item Average effort: 4--6 hours per day
		\item Training cadence: one course per week
		\item Initial readiness: approximately 9--10 weeks
		\item Professional mastery requires extended real-world practice
	\end{itemize}
	
	\newpage
	
	
	\newpage
	
	\section*{Part 2: Detailed Training Plan (Step-by-Step, Time-Based)}
	
	\subsection*{Assumptions}
	\begin{itemize}
		\item Each course takes 2–4 days
		\item Average effort: 4–6 hours/day
		\item Training pace: 1 course per week
		\item Total timeline below focuses on initial readiness (approx. 10–12 weeks)
		\item Full mastery continues beyond this with practice
	\end{itemize}
	
	\vspace{0.5cm}
	
	% ============================================
	% PHASE 1
	% ============================================
	\section*{ Phase 1: Foundation \& Working Style (Week 1)}
	
	\begin{tabularx}{\textwidth}{|X|p{3cm}|p{3cm}|X|}
		\hline
		\textbf{Week} & \textbf{Course} & \textbf{Time} & \textbf{Goal} \\
		\hline
		1 & \href{https://www.udemy.com/course/remote-independent-working-in-a-managerless-environment-101/}{Course 1: Remote \& Independent Working} & 2–3 days & Learn how to work independently; understand ticket‑based and manager‑less environments; set expectations for self‑driven technical work \\
		\hline
	\end{tabularx}
	
	\vspace{0.3cm}
	
	% ============================================
	% PHASE 2
	% ============================================
	\section*{ Phase 2: Infrastructure \& Core Python (Weeks 2–3)}
	
	\begin{tabularx}{\textwidth}{|X|p{3cm}|p{3cm}|X|}
		\hline
		\textbf{Week} & \textbf{Course} & \textbf{Time} & \textbf{Goal} \\
		\hline
		2 & \href{https://www.udemy.com/course/big-data-infra-admin-101-ticket-based-learning/}{Course 2: Big Data Infrastructure \& Ticket‑Based Learning} & 3–4 days & Understand big‑data systems; learn production‑style workflows; exposure to infrastructure issues and troubleshooting \\
		\hline
		3 & \href{https://www.udemy.com/course/python-full-stack-computational-framework-for-begineers-101/}{Course 3: Python Full‑Stack Computational Framework (Beginner)} & 3–4 days & Python fundamentals; data structures and computational thinking; writing clean, reusable Python code \\
		\hline
	\end{tabularx}
	
	\vspace{0.3cm}
	
	% ============================================
	% PHASE 3
	% ============================================
	\section*{ Phase 3: Testing, Debugging \& Intermediate Python (Weeks 4–6)}
	
	\begin{tabularx}{\textwidth}{|X|p{3cm}|p{3cm}|X|}
		\hline
		\textbf{Week} & \textbf{Course} & \textbf{Time} & \textbf{Goal} \\
		\hline
		4 & \href{https://www.udemy.com/course/full-spectrum-comprehensive-testing-error-reading-101/}{Course 4: Comprehensive Testing \& Error Reading} & 2–3 days & Debugging skills; understanding stack traces; reading and fixing errors independently \\
		\hline
		5 & \href{https://www.udemy.com/course/refresher-python-big-data-intermediate-quiz-based-course-102/}{Course 5: Python Big Data Refresher (Intermediate)} & 3–4 days & Reinforce Python concepts; apply Python to data‑heavy problems; build confidence in writing intermediate‑level code \\
		\hline
		6 & \href{https://www.udemy.com/course/python-full-stack-and-backend-engines-for-mc-ml-engines-102/}{Course 6: Python Backend \& ML Engines} & 3–4 days & Backend computation; understanding Python in ML / analytics pipelines; performance‑oriented programming \\
		\hline
	\end{tabularx}
	
	\vspace{0.3cm}
	
	% ============================================
	% PHASE 4
	% ============================================
	\section*{Phase 4: Testing at Scale \& Distributed Thinking (Weeks 7–9)}
	
	\begin{tabularx}{\textwidth}{|X|p{3cm}|p{3cm}|X|}
		\hline
		\textbf{Week} & \textbf{Course} & \textbf{Time} & \textbf{Goal} \\
		\hline
		7 & \href{https://www.udemy.com/course/writing-tests-for-simeng-python-code-conversion-concepts-101/}{Course 7: Writing Tests for Simulation \& Code Conversion} & 2–3 days & Writing tests for complex logic; ensuring correctness during refactoring; quality control mindset \\
		\hline
		8 & \href{https://www.udemy.com/course/full-stack-computational-framework-for-big-data-simulations/}{Course 8: Full‑Stack Computational Framework for Big Data} & 3–4 days & End‑to‑end computation pipelines; simulation‑based design; scaling logic from small to large systems \\
		\hline
		9 & \href{https://www.udemy.com/course/converting-standalone-code-to-distributed-code-for-analytics/}{Course 9: Converting Standalone Code to Distributed Code} & 3–4 days & Distributed computing concepts; parallel execution; production‑grade analytics workflows \\
		\hline
	\end{tabularx}
	
	\newpage
	
	% ============================================
	% SUMMARY TIMELINE
	% ============================================
	\section*{ Summary Timeline}
	
	\begin{tabularx}{\textwidth}{|l|X|}
		\hline
		\textbf{Phase} & \textbf{Duration} \\
		\hline
		Fundamentals \& Workflow & 1 week \\
		\hline
		Core Python \& Infrastructure & 2 weeks \\
		\hline
		Testing \& Intermediate Python & 3 weeks \\
		\hline
		Distributed Systems \& Scale & 3 weeks \\
		\hline
		\textbf{Initial Readiness} & \textbf{\~9–10 weeks} \\
		\hline
		Professional Mastery & 6–24 months with practice \\
		\hline
	\end{tabularx}
	
	\vspace{0.5cm}
	
	\section*{How to Use This Plan}
	\begin{itemize}
		\item Follow strict order
		\item Do hands‑on coding alongside videos
		\item Re‑do exercises if unclear
		\item Expect difficulty — that’s normal
		\item This is career‑level training, not casual learning
	\end{itemize}
	
	\newpage
	
	% ============================================
	% EXPANDED ACRONYMS AND PEDAGOGICAL BASIS
	% ============================================
	\section*{ Appendix: Expanded Acronyms \& Pedagogical Basis for the Course Sequence}
	
	\subsection*{Expanded Acronyms}
	\begin{itemize}
		\item \textbf{ML} — Machine Learning
		\item \textbf{MC} — Monte Carlo (simulation methods)
		\item \textbf{SimEng} — Simulation Engineering
		\item \textbf{Infra} — Infrastructure
		\item \textbf{Admin} — Administration
	\end{itemize}
	
	\subsection*{Why This Strict Order? (Pedagogical Explanation)}
	
	The training plan follows a \textbf{cumulative skill ladder}. Each course provides concepts and tools required for the next. Skipping or reordering breaks the scaffold.
	
	\begin{enumerate}
		\item \textbf{Week 1 (Remote Work)} — Before any technical skill, you learn \emph{how to learn independently} in a manager‑less, ticket‑driven environment. This sets the metacognitive foundation.
		
		\item \textbf{Week 2 (Big Data Infrastructure)} — You need to understand the \emph{ecosystem} (clusters, data pipelines, production issues) before writing code for it. Theory before practice.
		
		\item \textbf{Week 3 (Python Beginner)} — Core programming fundamentals. Without clean, reusable Python, later testing and distributed code become impossible.
		
		\item \textbf{Week 4 (Testing \& Error Reading)} — Debugging is taught \emph{early} so you can self‑correct. This prevents frustration in later, more complex courses.
		
		\item \textbf{Week 5 (Python Intermediate Refresher)} — Reinforcement through quiz‑based, data‑heavy problems solidifies Week 3 concepts before moving to backend work.
		
		\item \textbf{Week 6 (Backend \& ML Engines)} — Performance‑oriented Python for real pipelines. Builds directly on Week 5.
		
		\item \textbf{Week 7 (Writing Tests)} — Now you can write correct, refactorable code. Tests are non‑negotiable for distributed systems.
		
		\item \textbf{Week 8 (Full‑Stack Big Data Framework)} — End‑to‑end pipelines. Integrates everything from Weeks 2–7.
		
		\item \textbf{Week 9 (Standalone → Distributed)} — The capstone. You convert working single‑machine code to parallel, distributed execution. This is the core skill for the bridge role.
	\end{enumerate}
	
	\subsection*{Cognitive Load Management}
	\begin{itemize}
		\item Early weeks focus on \emph{one new major concept} per week.
		\item Testing and debugging are introduced \emph{before} complexity spikes (Week 4).
		\item Distributed computing is saved for \emph{last} (Week 9) because it requires all previous skills.
	\end{itemize}
	
	\subsection*{Time-to-Competence Realism}
	\begin{itemize}
		\item \textbf{9–10 weeks} → Initial readiness (can complete supervised tickets).
		\item \textbf{6 months} → Consistent independent work on medium‑complexity tasks.
		\item \textbf{1–2 years} → Full mastery (architecting distributed solutions, mentoring others).
	\end{itemize}
	
	\vspace{1cm}
	\noindent\textit{Document generated from the original training plan by Satyadhar. All Udemy links are preserved as provided.}
	
\end{document}